\section{Appointed texts and source locators}
The biblical English below is the public-domain Douay--Rheims/Challoner study translation, as retained from Project Gutenberg eBook 1581. It is not the approved English proclaimed at this celebration; the U.S. Lectionary controls that wording. Prayer and composed-antiphon bodies are withheld under the source rights recorded in the scope appendix.

\subsection{Entrance antiphon}
\label{proper-entrance-antiphon}
\textit{Gaudeamus omnes in Domino}. Complete text: \work{Missale Romanum}, All Saints, Entrance; English publisher witness, \href{https://www.liturgyoffice.org.uk/Missal/Music/Antiphonary.pdf}{ICEL Antiphonary}, p. 134 (PDF 142). No full English body is supplied.

\subsection{Collect}
\label{proper-collect}
\textit{Omnipotens sempiterne Deus}. Complete text: \work{Roman Missal}, Proper of Saints, 1 November, All Saints, Collect; qualified Latin witness, \href{https://media.musicasacra.com/books/latin_missal2002.pdf#page=512}{2002 secondary reproduction, PDF 512}. The current U.S. approved English body is not supplied.

\subsection{First Reading}
\label{proper-first-reading}
\properlocus{Revelation 7:2--4,9--14; Douay--Rheims title: Apocalypse.}
\begin{quote}
And I saw another angel ascending from the rising of the sun, having the sign of the living God. And he cried with a loud voice to the four angels to whom it was given to hurt the earth and the sea,

Saying: Hurt not the earth nor the sea nor the trees, till we sign the servants of our God in their foreheads.

And I heard the number of them that were signed. An hundred forty-four thousand were signed, of every tribe of the children of Israel.
\end{quote}
The appointment passes over verses 5--8.
\begin{quote}
After this, I saw a great multitude, which no man could number, of all nations and tribes and peoples and tongues, standing before the throne and in sight of the Lamb, clothed with white robes, and palms in their hands.

And they cried with a loud voice, saying: Salvation to our God, who sitteth upon the throne and to the Lamb.

And all the angels stood round about the throne and the ancients and the four living creatures. And they fell down before the throne upon their faces and adored God,

Saying: Amen. Benediction and glory and wisdom and thanksgiving, honour and power and strength, to our God, for ever and ever. Amen.

And one of the ancients answered and said to me: These that are clothed in white robes, who are they? And whence came they?

And I said to him: My Lord, thou knowest. And he said to me: These are they who are come out of great tribulation and have washed their robes and have made them white in the blood of the Lamb.
\end{quote}

\subsection{Responsorial Psalm}
\label{proper-responsorial-psalm}
\properlocus{Hebrew Psalm 24:1bc--2,3--4ab,5--6; Douay--Rheims/Vulgate Psalm 23.}
\begin{quote}
The earth is the Lord’s and the fulness thereof: the world, and all they that dwell therein.

For he hath founded it upon the seas; and hath prepared it upon the rivers.

Who shall ascend into the mountain of the Lord: or who shall stand in his holy place?

The innocent in hands, and clean of heart, who hath not taken his soul in vain.

He shall receive a blessing from the Lord, and mercy from God his Saviour.

This is the generation of them that seek him, of them that seek the face of the God of Jacob.
\end{quote}
The superscription and the final oath clause of verse 4 are omitted. The response is appointed as cf. verse 6. Its scriptural basis is the last sentence above; the approved adapted refrain is not supplied or reconstructed.

\subsection{Second Reading}
\label{proper-second-reading}
\properlocus{1 John 3:1--3.}
\begin{quote}
Behold what manner of charity the Father hath bestowed upon us, that we should be called and should be the sons of God. Therefore the world knoweth not us, because it knew not him.

Dearly beloved, we are now the sons of God: and it hath not yet appeared what we shall be. We know that when he shall appear we shall be like to him: because we shall see him as he is.

And every one that hath this hope in him sanctifieth himself, as he also is holy.
\end{quote}

\subsection{Gospel acclamation}
\label{proper-gospel-acclamation}
\properlocus{Matthew 11:28; scriptural basis of the appointed sung verse.}
\begin{quote}
Come to me all you that labor and are burdened, and I will refresh you.
\end{quote}
This is the study Bible verse, not a replacement for the approved Alleluia text and its attribution.

\subsection{Gospel}
\label{proper-gospel}
\properlocus{Matthew 5:1--12a.}
The study witness places meekness before mourning in verses 4--5; the U.S. proclaimed text has the reverse order. The final sentence about the prophets is beyond the appointed endpoint.
\begin{quote}
And seeing the multitudes, he went up into a mountain, and when he was set down, his disciples came unto him.

And opening his mouth he taught them, saying:

Blessed are the poor in spirit: for theirs is the kingdom of heaven.

Blessed are the meek: for they shall possess the land.

Blessed are they that mourn: for they shall be comforted.

Blessed are they that hunger and thirst after justice: for they shall have their fill.

Blessed are the merciful: for they shall obtain mercy.

Blessed are the clean of heart: for they shall see God.

Blessed are the peacemakers: for they shall be called the children of God.

Blessed are they that suffer persecution for justice’ sake: for theirs is the kingdom of heaven.

Blessed are ye when they shall revile you, and persecute you, and speak all that is evil against you, untruly, for my sake:

Be glad and rejoice for your reward is very great in heaven.
\end{quote}

\subsection{Prayer over the Offerings}
\label{proper-prayer-over-offerings}
\textit{Grata tibi sint, Domine, munera}. Complete text: \work{Roman Missal}, All Saints, Prayer over the Offerings; \href{https://media.musicasacra.com/books/latin_missal2002.pdf#page=512}{qualified Latin witness, PDF 512--513}. The current U.S. approved English body is not supplied.

\subsection{Proper Preface}
\label{proper-proper-preface}
\work{The glory of Jerusalem, our mother}. Complete publisher English witness: \href{https://www.icelweb.org/musicfolder/73JerusalemD.pdf}{ICEL, All Saints Preface, pp. 1--2}. Its Latin title is \textit{De gloria matris nostrae Ierusalem}. No full English or Latin body is supplied.

\subsection{Communion antiphon}
\label{proper-communion-antiphon}
\properlocus{Matthew 5:8--10; public-domain study Scripture, not approved Missal English.}
\begin{quote}
Blessed are the clean of heart: for they shall see God.

Blessed are the peacemakers: for they shall be called the children of God.

Blessed are they that suffer persecution for justice’ sake: for theirs is the kingdom of heaven.
\end{quote}
Complete approved antiphon: ICEL Antiphonary, All Saints, Communion, p. 134 (PDF 142).

\subsection{Prayer after Communion}
\label{proper-prayer-after-communion}
\textit{Mirabilem te, Deus}. Complete text: \work{Roman Missal}, All Saints, Prayer after Communion; \href{https://media.musicasacra.com/books/latin_missal2002.pdf#page=513}{qualified Latin witness, PDF 513--514}. The current U.S. approved English body is not supplied.

\subsection{Optional solemn blessing}
\label{proper-solemn-blessing}
All Saints, Solemn Blessing no. 18, \textit{De omnibus Sanctis}, meaning ``For All Saints.'' Complete text: \work{Roman Missal}, Solemn Blessings, no. 18; \href{https://media.musicasacra.com/books/latin_missal2002.pdf#page=377}{qualified Latin witness, PDF 377--378}. This blessing is permitted if selected. No full approved English body is supplied.

